%% 53장 난수 생성
\chapter{난수 생성}
\label{ch:random-numbers}

\abstract{몬테카를로 방법과 붓스트랩, 무작위 초기화는 모두 49장의 독립인 균등 확률변수의 무한열을 필요로 한다. 컴퓨터는 그것을 결정적인 계산으로 흉내 낸다. 이 장에서는 그 흉내가 어떻게 이루어지고 어디서 실패하는지를 다룬다. 선형 합동 생성기의 연속된 세 값이 공간의 몇 안 되는 평면 위에 놓인다는 Marsaglia의 발견을 RANDU에서 확인하고, 상태를 따라 걷는 생성기와 원소 번호에서 바로 난수를 만드는 계수기형 생성기를 견준다. stardust가 쓰는 Philox4×64-10의 라운드를 적고, 계수기형 설계가 스레드 수와 무관하게 같은 난수를 주는 까닭을 설명한다. 이어서 균등 정수를 치우침 없이 만드는 Lemire의 방법을 증명하고, Fisher--Yates 섞기와 지수 키를 쓴 가중치 비복원 추출이 정확한 분포를 준다는 것을 보인다. 끝으로 분포별 표본의 세 가지 기법, 곧 역변환, 기각, ziggurat를 다룬다.}

\section{의사난수}
\label{sec:pseudo-random}

결정적인 알고리즘이 만드는 수열이 정말 무작위일 수는 없다. 의사난수 생성기에 요구하는 것은 통계 검정으로 무작위 수열과 구별되지 않는다는 것이다. 생성기는 상태 $s$와 갱신 함수 $f$, 출력 함수 $g$로 이루어져 $s_{k+1}=f(s_k)$, $u_k=g(s_k)$를 만든다. 상태가 유한하므로 수열은 언젠가 되풀이되고, 그 \term{주기}{period}는 쓸 만큼보다 훨씬 길어야 한다. 또 연속된 $t$개 값의 순서쌍 $(u_k,\dots,u_{k+t-1})$이 $[0,1)^t$에 고르게 퍼져야 하는데, 이 다차원 균등성이 생성기의 질을 가장 날카롭게 가른다.

\subsection{선형 합동 생성기}

가장 오래된 생성기는 Lehmer의 \term{선형 합동 생성기}{linear congruential generator}
\[
x_{k+1}=(ax_k+c)\bmod m,\qquad u_k=x_k/m
\]
이다\cite{lehmer1951}. 계산이 곱셈 하나와 덧셈 하나로 빠르고, 주기와 같은 성질이 정수론으로 분석된다\cite{knuth1997taocp2}. 그러나 Marsaglia는 1968년에 이 생성기의 연속된 $t$개 값이 $[0,1)^t$에서 많아야 $(t!\,m)^{1/t}$개의 평행한 초평면 위에 놓인다는 것을 보였다\cite{marsaglia1968}. 1960년대에 IBM이 보급한 RANDU($a=65539$, $c=0$, $m=2^{31}$)가 극단적인 예이다. $a=2^{16}+3$이므로 $a^2=6a-9\pmod{2^{31}}$이고
\[
x_{k+2}=6x_{k+1}-9x_k\pmod{2^{31}}
\]
이다. 따라서 $9u_k-6u_{k+1}+u_{k+2}$는 정수이고, 그 값은 $-6<9u_k-6u_{k+1}+u_{k+2}<10$이어서 $-5$부터 $9$까지 $15$개뿐이다. 연속된 세 값으로 만든 점은 모두 $15$개의 평면 위에 놓인다. 실제로 시드 $1$에서 $30{,}000$개의 삼중쌍을 만들어 보면 정확히 $15$개의 평면이 나타난다. 3차원 몬테카를로 적분에 RANDU를 쓰면 공간의 대부분을 건너뛰게 된다.

\subsection{현대의 생성기}

이후의 생성기들은 상태를 키우고 갱신을 비선형으로 만들어 이 문제를 줄였다. Mersenne Twister는 $2^{19937}-1$의 주기와 $623$차원까지의 균등 분포를 보장하고\cite{matsumoto1998}, PCG는 작은 선형 합동 생성기의 출력에 비트 회전과 같은 순열을 씌워 격자 구조를 감춘다\cite{oneill2014pcg}. numpy의 기본 생성기가 PCG64이다. 생성기의 질은 TestU01의 BigCrush처럼 수백 개의 통계 검정으로 이루어진 묶음으로 잰다\cite{lecuyer2007}.

이 생성기들은 모두 상태를 따라 한 걸음씩 걷는다. $k$번째 값을 얻으려면 앞의 값들을 거쳐야 하고(점프 기능이 있어도 비싸다), 병렬로 쓰려면 스레드마다 다른 스트림을 나누어 주어야 한다. 그러면 결과가 스레드 수와 작업 분할에 의존한다.

\section{계수기형 생성기}
\label{sec:counter-based}

Salmon 등은 2011년에 다른 설계를 내놓았다\cite{salmon2011random123}. 상태를 갱신하는 대신, 키 $k$와 계수기 $c$를 받아 출력을 만드는 함수 $u=b_k(c)$를 두는 것이다. $c=0,1,2,\dots$에 차례로 적용하면 난수열이 되고, 어느 위치의 값이든 앞의 값 없이 바로 계산된다. 좋은 $b_k$는 암호학의 블록 암호와 같은 성질, 곧 입력의 비트 하나가 바뀌면 출력의 비트가 절반쯤 바뀌는 확산을 가져야 한다. 다만 암호학적 안전성은 필요 없으므로 라운드 수를 줄여 빠르게 만든다.

\term{Philox}{Philox}4×64는 64비트 단어 네 개의 계수기 $(c_0,c_1,c_2,c_3)$와 64비트 단어 두 개의 키 $(k_0,k_1)$에서 64비트 단어 네 개를 만든다. 한 라운드는
\begin{gather*}
(h_0,l_0)=M_0c_0,\qquad(h_1,l_1)=M_1c_2,\\
(c_0,c_1,c_2,c_3)\gets(h_1\oplus c_1\oplus k_0,\ l_1,\ h_0\oplus c_3\oplus k_1,\ l_0)
\end{gather*}
이다. 여기서 $(h,l)$은 $128$비트 곱의 위와 아래 $64$비트이고, 곱하는 상수와 라운드 사이에 키에 더하는 Weyl 상수는
\begin{gather*}
M_0=\texttt{0xD2E7470EE14C6C93},\qquad M_1=\texttt{0xCA5A826395121157},\\
\lfloor2^{64}/\varphi\rfloor=\texttt{0x9E3779B97F4A7C15},\qquad\lfloor2^{64}(\sqrt3-1)\rfloor=\texttt{0xBB67AE8584CAA73B}
\end{gather*}
이다. 뒤의 두 상수를 각각 $k_0$과 $k_1$에 라운드마다 더하고, 라운드를 열 번 되풀이한 것이 Philox4×64-10이다. 곱셈의 위 비트는 아래 비트들을 모두 섞고, 아래 비트는 다음 라운드의 위 비트로 옮겨 가므로, 열 라운드면 BigCrush를 통과하고도 여유가 있다.

계수기형 설계의 실용적인 장점은 결정성이다. stardust는 원소 $e$에 계수기 $(\lfloor e/4\rfloor,0,0,0)$ 블록의 $e\bmod4$번째 단어를 준다. 그러면 원소 하나의 값은 키와 $e$만으로 정해지므로, 배열을 몇 개의 스레드가 어떻게 나누어 채우든, 호출을 나누어 부르든 같은 값이 나온다. 기각 표본처럼 원소 하나가 단어를 더 필요로 하면 계수기 $(e,j,0,0)$($j\ge1$) 블록을 쓴다. 둘째 단어가 $0$인 블록과 겹치지 않으므로 원소들의 단어가 서로 섞이지 않는다.

\section{균등한 값}
\label{sec:uniform-values}

\subsection{실수}

64비트 단어 $w$에서 $[0,1)$의 실수를 만드는 표준 방법은 위 $53$비트를 취해 $2^{-53}$을 곱하는 것이다. 결과는 $\{0,2^{-53},\dots,1-2^{-53}\}$에서 균등하다. 이 값은 $0$이 될 수 있으므로 $-\log u$나 $\Phi^{-1}(u)$에 넣을 때는 조심해야 한다. 역변환 표본에서는 대신 $(m+1/2)2^{-53}$을 쓰면 $0$과 $1$이 나오지 않고 $0$과 $1$에 대칭이다.

\subsection{정수}

$[0,r)$의 균등 정수를 만드는 흔한 방법 $w\bmod r$은 $r$이 $2^{64}$을 나누지 않으면 치우친다. $2^{64}=qr+s$($0\le s<r$)이면 $0$부터 $s-1$까지의 값은 $q+1$번, 나머지는 $q$번 나오기 때문이다. 극단적인 예로 $r=3\cdot2^{62}$이면 $s=2^{62}$이어서 아래 $1/3$의 값이 나머지의 두 배 확률로 나온다. 치우침을 없애려면 $s$개의 단어를 버려야 하는데, Lemire는 나눗셈 없이 이것을 하는 방법을 보였다\cite{lemire2019interval}.

\begin{proposition}[Lemire]
\label{prop:lemire}
$w$가 $\{0,\dots,2^{64}-1\}$에서 균등하고 $s=2^{64}\bmod r$이라 하자. $128$비트 곱 $wr$의 아래 $64$비트 $\ell=wr\bmod2^{64}$이 $s$ 이상일 때만 받아들이면, 받아들인 $w$의 위 $64$비트 $\lfloor wr/2^{64}\rfloor$는 $\{0,\dots,r-1\}$에서 균등하다.
\end{proposition}

\begin{proof}
값 $j\in\{0,\dots,r-1\}$에 대해 $\lfloor wr/2^{64}\rfloor=j$인 $w$는 $wr\in[j2^{64},(j+1)2^{64})$인 것들이다. 이 구간에서 $r$의 배수 $wr$의 아래 $64$비트 $\ell=wr-j2^{64}$은 $\ell\equiv-j2^{64}\pmod r$인 $[0,2^{64})$의 수들이다. 그런 수 가운데 $[s,2^{64})$에 드는 것의 개수를 세자. $[s,2^{64})$의 길이는 $2^{64}-s=qr$로 $r$의 배수이므로, 그 안에는 법 $r$의 각 잉여류가 정확히 $q$개씩 있다. 따라서 모든 $j$에 대해 받아들여지는 $w$가 정확히 $q$개이고, 위 $64$비트는 균등하다.
\end{proof}

받아들이지 않을 확률은 $s/2^{64}<r/2^{64}$로 작은 $r$에서는 무시할 만하다. 또 $\ell\ge r$이면 $\ell\ge s$가 저절로 성립하므로, 대부분의 경우 $s$를 계산하는 나눗셈조차 필요 없다. $\ell<r$일 때만 $s=(2^{64}-r)\bmod r$을 계산해 판정한다.

\section{섞기와 추출}
\label{sec:shuffle}

\subsection{Fisher--Yates}

길이 $n$인 배열을 균등한 무작위 순열로 섞는 표준 방법은 $i=0,1,\dots,n-2$에 대해 $[i,n)$의 균등 정수 $j$를 뽑아 $a_i$와 $a_j$를 바꾸는 것이다. Fisher와 Yates가 1938년에 손으로 하는 방법으로 적었고\cite{fisher1938tables}, Durstenfeld가 지금의 제자리 판으로 적었다\cite{durstenfeld1964}.

\begin{proposition}
Fisher--Yates 섞기는 $n!$개의 순열을 모두 확률 $1/n!$로 낸다.
\end{proposition}

\begin{proof}
걸음 $i$에서 $a_i$의 자리에 오는 것은 아직 놓이지 않은 $n-i$개 원소 가운데 균등하게 고른 하나이다. 걸음들의 선택은 독립이므로 특정한 순열이 나올 확률은 $\prod_{i=0}^{n-2}1/(n-i)=1/n!$이다. 서로 다른 선택열은 서로 다른 순열을 주므로 모든 순열이 이 확률로 나온다.
\end{proof}

자주 보는 잘못은 $j$를 $[i,n)$ 대신 $[0,n)$에서 뽑는 것이다. 그러면 선택열이 $n^n$개인데 $n^n$은 $n!$의 배수가 아니므로($n\ge3$) 균등할 수 없다. 앞의 $k$걸음만 하면 크기 $k$의 균등한 비복원 표본이 된다.

\subsection{가중치 비복원 추출}

원소마다 무게 $p_i$가 있을 때, 하나를 무게에 비례해 뽑고 그것을 빼고 다시 뽑는 일을 $k$번 하는 것이 가중치 비복원 추출이다. Efraimidis와 Spirakis는 이것을 한 번의 정렬로 하는 방법을 보였다\cite{efraimidis2006}. 독립인 지수 확률변수 $E_i$로 키 $E_i/p_i$를 만들고 가장 작은 $k$개를 고르면 된다. $E_i/p_i$는 비율 $p_i$인 지수 분포를 따르고, 독립인 지수 확률변수들 가운데 최솟값이 $i$일 확률이 $p_i/\sum_jp_j$이며, 지수 분포의 무기억성으로 남은 원소들의 키는 최솟값을 뺀 뒤에도 같은 분포이다. 그래서 키의 순서가 정확히 한 번에 하나씩 무게에 비례해 뽑는 순서와 같은 분포를 가진다.

\section{분포별 표본}
\label{sec:variate-generation}

\subsection{역변환}

49장의 역변환 정리로 $F^{-}(U)$는 분포 $F$를 따른다. 분위수 함수를 정확히 계산할 수 있으면 이것이 가장 일반적인 방법이고, 단조 변환이므로 공통 난수 기법이나 준몬테카를로와 잘 어울린다. 지수 분포 $-\log(1-U)$, 코시 $\tan(\pi(U-1/2))$, 이산 분포의 순차 탐색이 그 예이다. 단점은 속도이다. 감마나 베타처럼 분위수 함수가 뉴턴 반복을 요구하는 분포에서는 표본 하나에 수백 나노초가 든다.

\subsection{기각}

\begin{proposition}[von Neumann]
\label{prop:rejection}
밀도 $f$와 $g$가 모든 $x$에서 $f(x)\le cg(x)$를 만족한다고 하자. $Y\sim g$와 독립인 $U\sim U(0,1)$을 뽑아 $U\le f(Y)/(cg(Y))$이면 $Y$를 받아들이고 아니면 다시 뽑는다. 받아들인 값은 $f$를 따르고, 한 번 시도에서 받아들일 확률은 $1/c$이다.
\end{proposition}

\begin{proof}
한 번의 시도에서 $P(Y\le x,\ \text{받아들임})=\int_{-\infty}^xg(y)\cdot f(y)/(cg(y))\,dy=F(x)/c$이다. $x\to\infty$로 보내면 받아들일 확률이 $1/c$이고, 시도들은 독립이므로 받아들인 값의 분포는 $P(Y\le x\mid\text{받아들임})=F(x)$이다.
\end{proof}

이 방법은 von Neumann이 1951년에 적었다\cite{vonneumann1951}. 시도 횟수는 평균 $c$번의 기하 분포이므로 $c$가 $1$에 가까운 덮개 $g$를 찾는 것이 기술이다. 예컨대 표준 정규를 코시 덮개로 뽑으면 $c=\sqrt{2\pi/e}=1.520$이어서 받아들임 확률이 $0.658$이다. Marsaglia와 Tsang의 감마 표본은 감마 확률변수를 정규 확률변수의 세제곱 변환으로 근사하고 기각으로 보정해, 모양이 $1$ 이상이면 받아들임 확률이 $0.95$를 넘는다\cite{marsaglia2000gamma}. 큰 평균의 Poisson과 이항에는 Hörmann의 변환 기각 방법(PTRS, BTRS)이 표준이다\cite{hormann1993poisson,hormann1993binomial}.

\subsection{Ziggurat}

정규와 지수처럼 단조 감소하는 꼬리를 가진 밀도에는 \term{ziggurat 방법}{ziggurat method}이 가장 빠르다\cite{marsaglia2000ziggurat}. 밀도 아래 영역을 넓이가 같은 가로 띠 $n$개(흔히 $128$이나 $256$)로 덮고, 맨 아래 띠는 꼬리를 포함하게 한다. 띠 하나를 균등하게 고르고 그 띠의 가로 범위에서 $x$를 균등하게 뽑았을 때, $x$가 위 띠의 오른쪽 끝보다 안쪽이면 그 점은 밀도 아래에 확실히 있으므로 곧바로 받아들인다. 이 빠른 경로는 무작위 단어 하나의 비트로 띠와 $x$를 함께 정하고, 표에서 읽은 곱셈과 비교 하나로 끝나며, 대부분의 시도가 여기서 끝난다. 띠의 가장자리에 떨어진 드문 경우에만 밀도를 계산해 기각 판정을 하고, 맨 아래 띠에서 꼬리로 나간 경우는 꼬리 전용 방법으로 뽑는다.

정규 표본에는 Box--Muller 변환 $\sqrt{-2\log U_1}(\cos2\pi U_2,\sin2\pi U_2)$도 있다\cite{box1958}. 극좌표로 바꾸면 독립 표준 정규 둘의 결합 밀도가 반지름의 제곱에 대해 지수 분포, 각도에 대해 균등 분포로 갈라진다는 데서 나온다. 정확하지만 로그와 삼각함수가 들어 ziggurat보다 느리다.

\begin{instardust}[\code{docs/design/random.md} §1--4, \code{docs/design/stats.md} §4]
\code{sd.random}의 생성기는 Philox4×64-10이며, Random123의 알려진 답 벡터와 비트까지 같다. 시드는 numpy의 \code{SeedSequence}와 같은 해시로 키를 만들지만, 계수기를 쓰는 방법이 달라 값은 numpy와 다르다. 이것은 numpy와의 비트 호환 대신 스레드 수와 무관한 결정성을 고른 결과이다. 한 블록은 곱셈 스무 번의 의존 사슬이어서 지연 시간에 묶이므로, 독립인 블록 넷의 라운드를 섞어 곱셈기를 채운다. \code{integers}는 Lemire의 방법이고, 범위가 하나인 호출은 $s$를 한 번만 계산한다. \code{permutation}은 Fisher--Yates이고, 비복원 균등 \code{choice}는 그 앞 $k$걸음을 건드린 항목만 담은 해시 표로 한다. 가중치 비복원 \code{choice}는 지수 키 $E/p$의 가장 작은 $k$개이다. 분포의 \code{sample}은 31개 모두 방법을 재서 골랐다. 정규와 로그정규는 256층 ziggurat, 지수 계열은 지수의 ziggurat와 닫힌 꼴 변환, 감마는 Marsaglia--Tsang, Poisson과 이항은 평균이 $10$ 이상이면 PTRS와 BTRS이고 그 아래는 역변환이다. 기각 판정의 로그 확률은 double-double로 계산한다. 평균이 $10^{12}$이면 로그 감마 함수의 차가 $10^{13}$ 크기여서 binary64로는 판정이 틀리기 때문이다.
\end{instardust}

\section{맺음}

의사난수 생성기는 결정적인 계산으로 독립인 균등 확률변수를 흉내 낸다. 그 흉내의 질은 다차원 균등성에서 드러나고, 선형 합동 생성기의 평면 구조가 그 실패의 고전적인 예이다. 계수기형 생성기는 상태를 없애 원소 번호에서 바로 난수를 만들므로, 병렬 계산에서도 스레드 수와 무관하게 같은 결과를 준다. 균등한 값을 만드는 단계에도 치우침의 함정이 있어서, 정수는 Lemire의 기각으로, 순열은 Fisher--Yates로 정확한 균등성을 얻는다. 분포별 표본은 역변환, 기각, ziggurat 가운데 분포마다 가장 빠르면서 정확한 것을 고른다. 다음 장부터는 이렇게 만든 표본과 실제 자료로 모수를 추정하는 통계적 추론으로 넘어간다.

난수 생성의 이론은 Knuth의 책\cite{knuth1997taocp2} 3장에, 분포별 표본의 방법은 Devroye의 책\cite{devroye1986}에 망라되어 있다.

\chapterreferences
